\documentclass[letterpaper]{article}
\usepackage{aaai24}
\usepackage{times}
\usepackage{helvet}
\usepackage{courier}
\usepackage[hyphens]{url}
\usepackage{graphicx}
\urlstyle{rm}
\def\UrlFont{\rm}
\usepackage{natbib}
\usepackage{caption}
\usepackage{amsmath}
\usepackage{booktabs}
\usepackage{listings}
\usepackage{placeins}
\frenchspacing
\setlength{\pdfpagewidth}{8.5in}
\setlength{\pdfpageheight}{11in}
\pdfinfo{/TemplateVersion (2024.1)}
\setcounter{secnumdepth}{1}
\lstset{basicstyle=\footnotesize\ttfamily,breaklines=true,showstringspaces=false}

\title{Optimization and Regularization in Dog Breed Classification}
\author{\input{generated/author}}
\affiliations{AI6103 Homework Assignment}

\begin{document}
\maketitle
\input{generated/status}

\begin{abstract}
\input{generated/abstract}
\end{abstract}

\section{Experimental Protocol}
We train torchvision EfficientNetV2-S with random initialization and a
120-class output layer \citep{tan2021efficientnetv2}. All experiments use
cross-entropy loss and SGD with momentum 0.9 \citep{goodfellow2016deep},
with batch size 128 and seed 42.
Each condition is trained independently.

Three learning rates are compared for 15 epochs each. The rate with the
lowest final training loss is then used for the 300-epoch schedule,
weight-decay, and mixup experiments. Weight decay and the final model are
selected by final validation accuracy, with validation loss breaking ties.
All results are measured at the final epoch; the selected model is evaluated
once on the test set after the experiments.

\section{Data Preprocessing}
\input{generated/data_summary}

RGB images are resized using bilinear interpolation to $100\times100$.
For each channel $c$, over the $N=\TrainingCount\times100\times100$ resized training
pixels scaled into $[0,1]$, we compute
\begin{equation}
 \mu_c=\frac{1}{N}\sum_{i=1}^{N}x_{ic},\qquad
 \sigma_c=\sqrt{\frac{1}{N}\sum_{i=1}^{N}(x_{ic}-\mu_c)^2}.
\end{equation}
The normalized input is $z_{ic}=(x_{ic}-\mu_c)/\sigma_c$.
The same training-set statistics normalize validation and test images.
\input{generated/statistics}

Training uses random horizontal flipping and random cropping with 12-pixel
zero padding, followed by normalization \citep{goodfellow2016deep}:
\begin{lstlisting}[language=Python]
T.Compose([
    T.Resize((100, 100), antialias=True),
    T.RandomHorizontalFlip(p=0.5),
    T.RandomCrop(100, padding=12,
        padding_mode="constant", fill=0),
    T.ToTensor(),
    T.Normalize(mean, std),
])
\end{lstlisting}
Validation and test use only resize, tensor conversion, and normalization.
Training loss and accuracy are evaluated on unaugmented images after each
epoch, using the same evaluation mode as validation. Loss denotes
cross-entropy without the weight penalty throughout the report.

\input{generated/results_table}

\section{Initial Learning Rate}
We compare $\eta_0\in\{0.5,0.005,0.001\}$ for 15 epochs, with no weight
decay or schedule. Random crop and horizontal flip are enabled in all runs.
\input{generated/lr_observations}
\IfFileExists{../artifacts/figures/learning_rate_report.pdf}{%
\begin{figure}[t]
\centering
\includegraphics[width=\columnwidth]{../artifacts/figures/learning_rate_report.pdf}
\caption{Initial-rate experiments: training and validation loss (top) and
top-1 accuracy (bottom) against epoch. Colors identify conditions; solid and dashed lines denote training and validation.}
\label{fig:lr}
\end{figure}}{}

\section{Learning Rate Schedule}
Cosine annealing reduces the rate smoothly, first slowly, then more rapidly,
and finally slowly as it approaches zero. Without warm restarts, its rule is
\begin{equation}
 \eta_t=\frac{\eta_0}{2}\left(1+\cos\frac{\pi t}{T}\right),
 \quad T=300,\quad 0\leq t\leq T.
\end{equation}
Here $t$ is the number of completed epochs. The schedule decreases from
$\eta_0$ to zero over 300 epochs \citep{loshchilov2017sgdr}.

We compare constant and cosine rates at the selected $\eta_0$ for 300
epochs, with all other settings unchanged and weight decay zero.
\input{generated/schedule_observations}
\IfFileExists{../artifacts/figures/schedule_report.pdf}{%
\begin{figure}[t]
\centering
\includegraphics[width=\columnwidth]{../artifacts/figures/schedule_report.pdf}
\caption{Constant versus cosine schedule over 300 epochs. Solid: training; dashed: validation.}
\label{fig:schedule}
\end{figure}}{}
\section{Weight Decay}
The two prescribed coefficients are $\lambda=5\times10^{-4}$ and
$10^{-4}$, with cosine annealing and the selected initial learning rate.
Both runs last 300 epochs. The zero-decay cosine run is reused as the control.
Weight decay penalizes large parameters \citep{zhang2019mechanisms}.
With SGD and momentum, the update is
\begin{align}
 g_t &= \nabla_w L(w_t)+\lambda w_t,\\
 v_t &= 0.9v_{t-1}+g_t,\qquad w_{t+1}=w_t-\eta_t v_t.
\end{align}
Here $L$ is the cross-entropy loss, and $\lambda w_t$ shrinks the weights
during each update. The same coefficient is applied to all parameters.

\input{generated/decay_observations}
\IfFileExists{../artifacts/figures/weight_decay_report.pdf}{%
\begin{figure}[t]
\centering
\includegraphics[width=\columnwidth]{../artifacts/figures/weight_decay_report.pdf}
\caption{Cosine-scheduled control and both weight-decay coefficients.
The CE excludes the penalty. Solid: training; dashed: validation.}
\label{fig:decay}
\end{figure}}{}
\section{Mixup and Final Evaluation}
For a random within-batch pairing $(i,j)$, mixup constructs
\begin{align}
 \lambda &\sim \operatorname{Beta}(0.2,0.2),\\
 \widetilde{x} &= \lambda x_i+(1-\lambda)x_j,\\
 \widetilde{y} &= \lambda y_i+(1-\lambda)y_j.
\end{align}
The loss is $\lambda\operatorname{CE}(f(\widetilde{x}),y_i)
+(1-\lambda)\operatorname{CE}(f(\widetilde{x}),y_j)$
\citep{zhang2018mixup}. We draw one coefficient per batch and permute
examples within the batch after preprocessing.

The beta probability density is
\begin{equation}
 p(\lambda)=\frac{\lambda^{-0.8}(1-\lambda)^{-0.8}}{B(0.2,0.2)},
 \quad 0<\lambda<1.
\end{equation}
The density is U-shaped (Figure~\ref{fig:beta}), so most mixtures are
dominated by one of the two images.
\begin{figure}[t]
\centering
\includegraphics[width=\columnwidth]{../artifacts/figures/beta_alpha_0.2.pdf}
\caption{Probability density of $\operatorname{Beta}(0.2,0.2)$.}
\label{fig:beta}
\end{figure}

The mixup run lasts 300 epochs using the best pre-mixup configuration.
\input{generated/mixup_observations}
\IfFileExists{../artifacts/figures/mixup_report.pdf}{%
\begin{figure}[t]
\centering
\includegraphics[width=\columnwidth]{../artifacts/figures/mixup_report.pdf}
\caption{Selected control versus mixup. Solid: training; dashed: validation.}
\label{fig:mixup}
\end{figure}}{}
\input{generated/test}
\FloatBarrier
\label{body-end}
\bibliography{references}
\end{document}
